% Complete Section 6.1, original PDF 73-78 / printed 63-68, including all Table 6.1 cells.
% block: exterior-algebra-heading page:73 type:section
\subsection{\BT{Exterior algebra}{外代数}}\label{sec:exterior-algebra}

% block: exterior-algebra-introduction page:73 type:ordinary-paragraph
\BT{The primary objects of study in exterior calculus are differential $k$-forms, which we take up in the next section. These are functions whose values are alternating multilinear forms, also called algebraic $k$-forms. The exterior calculus is naturally concerned with the operations of differentiation and integration of these functions. But first we establish the basic terminology and properties of the algebraic $k$-forms which arise at each point, just as, before studying vector calculus, we must first understand finite dimensional vector spaces. This is the subject of exterior algebra, initiated by Grassman in the middle of the 19th century. Here, we briefly review the definitions and main properties, all of which are elementary.}{外微积分的主要研究对象是微分 $k$-形式, 下一节将讨论它们. 这些函数的值是交错多重线性形式, 也称为代数 $k$-形式. 外微积分自然涉及这些函数的微分和积分运算. 但首先需要建立每个点处出现的代数 $k$-形式的基本术语和性质, 正如在研究向量微积分之前, 必须先理解有限维向量空间一样. 这正是外代数的内容, 它由 Grassman 在 19 世纪中叶创立. 这里简要回顾其定义和主要性质, 这些都是初等的.}

% block: multilinear-forms-heading page:73 type:subsection
\subsubsection{\BT{Multilinear forms}{多重线性形式}}\label{subsec:multilinear-forms}

% block: multilinear-forms-definition pages:73,74 type:ordinary-paragraph
\BT{For a finite dimensional vector space $V$, and a nonnegative integer $k$, we consider the vector space $\operatorname{Lin}^kV$ consisting of all $k$-linear maps}{对有限维向量空间 $V$ 和非负整数 $k$, 考虑由所有 $k$-线性映射构成的向量空间 $\operatorname{Lin}^kV$:}
\[
\omega:\overbrace{V\times\cdots\times V}^{k\text{ times}}\to\mathbb R.
\]
\BT{By convention, we understand $\operatorname{Lin}^0V$ to be the space of scalars $\mathbb R$. If $\omega\in\operatorname{Lin}^jV$ and $\mu\in\operatorname{Lin}^kV$, the tensor product $\omega\otimes\mu\in\operatorname{Lin}^{j+k}V$ is given by}{按约定, 将 $\operatorname{Lin}^0V$ 理解为标量空间 $\mathbb R$. 若 $\omega\in\operatorname{Lin}^jV$, $\mu\in\operatorname{Lin}^kV$, 则张量积 $\omega\otimes\mu\in\operatorname{Lin}^{j+k}V$ 由下式给出:}
\[
\begin{gathered}
(\omega\otimes\mu)(v_1,\ldots,v_{j+k})=\omega(v_1,\ldots,v_j)\mu(v_{j+1},\ldots,v_{j+k}),\\
v_1,\ldots,v_{j+k}\in V.
\end{gathered}
\]

% block: multilinear-dual-basis page:74 type:ordinary-paragraph
\BT{From the definition, the space $\operatorname{Lin}^1V$ is just the dual space $V'$. We refer to its elements as either \emph{covectors} or $1$-forms. If we choose a basis $u_1,\ldots,u_n$ of $V$ ($n=\dim V$), this induces a dual basis $u^1,\ldots,u^n$ of $V'$, defined by $u^i(u_j)=\delta_{ij}$. A general element of $\operatorname{Lin}^kV$ is determined by the value it assigns to each of the $k$-tuples $(u_{\sigma_1},\ldots,u_{\sigma_k})$ with the $\sigma_i$ integers between $1$ and $n$. Thus $\dim\operatorname{Lin}^kV=n^k$. The element of $\operatorname{Lin}^kV$ which takes the $k$-tuple $(u_{\sigma_1},\ldots,u_{\sigma_k})$ to $1$ and all the other $k$-tuples $(u_{\tau_1},\ldots,u_{\tau_k})$ to $0$ is simply $u^{\sigma_1}\otimes\cdots\otimes u^{\sigma_k}$, and so the set of all of these gives a basis for $\operatorname{Lin}^kV$.}{由定义, 空间 $\operatorname{Lin}^1V$ 就是对偶空间 $V'$. 其元素称为\emph{余向量}或 $1$-形式. 若选取 $V$ 的一组基 $u_1,\ldots,u_n$ ($n=\dim V$), 就诱导了 $V'$ 的对偶基 $u^1,\ldots,u^n$, 由 $u^i(u_j)=\delta_{ij}$ 定义. $\operatorname{Lin}^kV$ 的一般元素由它赋予每个 $k$ 元组 $(u_{\sigma_1},\ldots,u_{\sigma_k})$ 的值决定, 其中 $\sigma_i$ 是 $1$ 到 $n$ 之间的整数. 因此 $\dim\operatorname{Lin}^kV=n^k$. 将 $k$ 元组 $(u_{\sigma_1},\ldots,u_{\sigma_k})$ 映为 $1$, 并将所有其他 $k$ 元组 $(u_{\tau_1},\ldots,u_{\tau_k})$ 映为 $0$ 的 $\operatorname{Lin}^kV$ 元素, 就是 $u^{\sigma_1}\otimes\cdots\otimes u^{\sigma_k}$, 因而所有这些元素构成 $\operatorname{Lin}^kV$ 的一组基.}

% block: multilinear-canonical-coordinate-basis page:74 type:ordinary-paragraph
\BT{The case $V=\mathbb R^n$ differs from the general case only in that there is a canonical basis $e_1,\ldots,e_n$. As notation for the corresponding dual basis we use $dx^1,\ldots,dx^n$. Thus, $dx^i$ is simply the linear functional which assigns to each vector $\mathbb R^n$ its $i$th component. The reason for the odd notation ($dx^i$ rather than, say, $e^i$) will become clear as we develop exterior calculus.}{情形 $V=\mathbb R^n$ 与一般情形的区别, 仅在于它具有典范基 $e_1,\ldots,e_n$. 相应对偶基记作 $dx^1,\ldots,dx^n$. 因此, $dx^i$ 就是将 $\mathbb R^n$ 中每个向量映为其第 $i$ 个分量的线性泛函. 随着外微积分的展开, 采用这一特别记号 ($dx^i$ 而不是例如 $e^i$) 的原因将变得清楚.}

% block: alternating-forms-heading page:74 type:subsection
\subsubsection{\BT{Alternating multilinear forms}{交错多重线性形式}}\label{subsec:alternating-multilinear-forms}

% block: alternating-forms-definition page:74 type:ordinary-paragraph
\BT{Multilinear forms $\omega\in\operatorname{Lin}^kV$ that change sign under the interchange of two variables are called \emph{skew-symmetric} or \emph{alternating}:}{交换两个变量时变号的多重线性形式 $\omega\in\operatorname{Lin}^kV$, 称为\emph{反对称}或\emph{交错}形式:}
\[
\begin{gathered}
\omega(v_1,\ldots,v_i,\ldots,v_j,\ldots,v_k)=-\omega(v_1,\ldots,v_j,\ldots,v_i,\ldots,v_k),\\
1\leq i<j\leq n,\qquad v_1,\ldots,v_k\in V.
\end{gathered}
\]
\BT{In particular, an alternating form vanishes if any of its arguments is repeated. To any multilinear form $\omega$ there is associated the alternating form which is its skew-symmetric part:}{特别地, 若交错形式的任一自变量重复出现, 则其值为零. 任意多重线性形式 $\omega$ 都伴随一个交错形式, 即其反对称部分:}
\[
(\operatorname{skw}\omega)(v_1,\ldots,v_k)=\frac1{k!}\sum_{\sigma\in\Sigma_k}\operatorname{sign}(\sigma)\omega(v_{\sigma_1},\ldots,v_{\sigma_k}),
\]
\BT{where the sum is over all the permutations of the integers $1$ to $k$. For example, if $\omega\in\operatorname{Lin}^3V$, then}{其中对整数 $1$ 到 $k$ 的所有排列求和. 例如, 若 $\omega\in\operatorname{Lin}^3V$, 则}
\[
\begin{aligned}
(\operatorname{skw}\omega)(v_1,v_2,v_3)=\frac16[&\omega(v_1,v_2,v_3)+\omega(v_2,v_3,v_1)+\omega(v_3,v_1,v_2)\\
&-\omega(v_1,v_3,v_2)-\omega(v_3,v_2,v_1)-\omega(v_2,v_1,v_3)].
\end{aligned}
\]

% block: alternating-space-dimension page:74 type:ordinary-paragraph
\BT{We define the space $\operatorname{Alt}^kV$ as the space of all skew-symmetric $k$-linear forms. The spaces $\operatorname{Alt}^0V$ and $\operatorname{Alt}^1V$ coincide with $\operatorname{Lin}^0V$ and $\operatorname{Lin}^1V$, i.e., are just the space of scalars and the space of covectors, respectively. An element of $\operatorname{Alt}^kV$ is determined by the value it assigns to each of the $k$-tuples $(e_{\sigma_1},\ldots,e_{\sigma_k})$ with $1\leq\sigma_1<\cdots<\sigma_k\leq n$, so $\dim\operatorname{Alt}^kV=\binom nk$. For $k>n$, $\operatorname{Alt}^kV=0$.}{将空间 $\operatorname{Alt}^kV$ 定义为所有反对称 $k$-线性形式构成的空间. 空间 $\operatorname{Alt}^0V$ 和 $\operatorname{Alt}^1V$ 分别与 $\operatorname{Lin}^0V$ 和 $\operatorname{Lin}^1V$ 一致, 即分别是标量空间和余向量空间. $\operatorname{Alt}^kV$ 的一个元素, 由它赋予满足 $1\leq\sigma_1<\cdots<\sigma_k\leq n$ 的每个 $k$ 元组 $(e_{\sigma_1},\ldots,e_{\sigma_k})$ 的值决定, 所以 $\dim\operatorname{Alt}^kV=\binom nk$. 当 $k>n$ 时, $\operatorname{Alt}^kV=0$.}

% block: exterior-product-definition pages:74,75 type:ordinary-paragraph
\BT{The \emph{exterior product} or \emph{wedge product} of $\omega\in\operatorname{Alt}^jV$ and $\mu\in\operatorname{Alt}^kV$ is given by}{$\omega\in\operatorname{Alt}^jV$ 和 $\mu\in\operatorname{Alt}^kV$ 的\emph{外积}或\emph{楔积}由下式给出:}
\[
(\omega\wedge\mu)(v_1,\ldots,v_{j+k})=\sum_\sigma\operatorname{sign}(\sigma)\omega(v_{\sigma_1},\ldots,v_{\sigma_j})\mu(v_{\sigma_{j+1}},\ldots,v_{\sigma_{j+k}}),
\]
\BT{where the sum is over $\sigma\in\Sigma_{j+k}$ for which $\sigma_1<\cdots<\sigma_j$ and $\sigma_{j+1}<\cdots<\sigma_{j+k}$. Equivalently,}{其中对满足 $\sigma_1<\cdots<\sigma_j$ 和 $\sigma_{j+1}<\cdots<\sigma_{j+k}$ 的 $\sigma\in\Sigma_{j+k}$ 求和. 等价地,}
\[
\omega\wedge\mu=\binom{j+k}j\operatorname{skw}(\omega\otimes\mu).
\]
\BT{The exterior product satisfies the anticommutativity law}{外积满足反交换律}
\[
\omega\wedge\mu=(-)^{jk}\mu\wedge\omega,\qquad\omega\in\operatorname{Alt}^jV,\ \mu\in\operatorname{Alt}^kV.
\]

% block: alternating-basis page:75 type:ordinary-paragraph
\BT{Let $\{u_i\}$ be a basis for $V$ and $\{u^i\}$ the corresponding dual basis. For given $1\leq\sigma_1<\cdots<\sigma_k\leq n$, the element of $\operatorname{Alt}^kV$ which takes the $k$-tuple $(u_{\sigma_1},\ldots,u_{\sigma_k})$ to $1$ and takes all the other $k$-tuples with increasing indices to $0$ is}{设 $\{u_i\}$ 是 $V$ 的一组基, $\{u^i\}$ 是相应的对偶基. 给定 $1\leq\sigma_1<\cdots<\sigma_k\leq n$, 将 $k$ 元组 $(u_{\sigma_1},\ldots,u_{\sigma_k})$ 映为 $1$, 将所有其他具有递增指标的 $k$ 元组映为 $0$ 的 $\operatorname{Alt}^kV$ 元素为}
\begin{equation}\label{eq:alternating-form-basis}
u^{\sigma_1}\wedge\cdots\wedge u^{\sigma_k}.
\end{equation}
\BT{The set of all of such form a basis for $\operatorname{Alt}^kV$, and so the general alternating $k$-form on $V$ is a linear combination of these. In particular, the general alternating $k$-form on $\mathbb R^n$ can be written}{所有这样的元素构成 $\operatorname{Alt}^kV$ 的一组基, 因而 $V$ 上的一般交错 $k$-形式是这些元素的线性组合. 特别地, $\mathbb R^n$ 上的一般交错 $k$-形式可写为}
\[
\sum_{1\leq\sigma_1<\cdots<\sigma_k\leq n}a_\sigma\,dx^{\sigma_1}\wedge\cdots\wedge dx^{\sigma_k}
\]
\BT{for a unique choice of coefficients $a_\sigma\in\mathbb R$.}{其中系数 $a_\sigma\in\mathbb R$ 唯一确定.}

% block: alternating-operations-heading page:75 type:subsection
\subsubsection{\BT{Operations on alternating forms}{交错形式上的运算}}\label{subsec:alternating-form-operations}

% block: alternating-contraction page:75 type:ordinary-paragraph
\medskip\noindent\textbf{\BT{Contraction}{缩并}}\par
\BT{Besides the vector space operations on $\operatorname{Alt}^kV$ and the exterior product, there are several other useful algebraic operations. Given $\omega\in\operatorname{Alt}^kV$ with $k\geq1$ and $v\in V$, the \emph{contraction} of $\omega$ and $v$, also known as the interior product, is the $(k-1)$-form}{除 $\operatorname{Alt}^kV$ 上的向量空间运算和外积外, 还有其他几种有用的代数运算. 给定 $\omega\in\operatorname{Alt}^kV$, 其中 $k\geq1$, 以及 $v\in V$, $\omega$ 与 $v$ 的\emph{缩并}, 也称为内乘, 是如下 $(k-1)$-形式:}
\[
(\omega\lrcorner v)(v_1,\ldots,v_{k-1})=\omega(v,v_1,\ldots,v_{k-1}).
\]
\BT{(If $k=0$, $\omega\lrcorner v:=0$.) Since $\omega$ is alternating, $(\omega\lrcorner v)\lrcorner v$ always vanishes. Moreover, it is easy to check that contraction with a vector is an antiderivation in the sense that}{(若 $k=0$, 规定 $\omega\lrcorner v:=0$.) 由于 $\omega$ 是交错的, $(\omega\lrcorner v)\lrcorner v$ 总为零. 此外, 容易验证, 与一个向量的缩并是一种反导子, 即}
\[
\begin{gathered}
(\omega\wedge\eta)\lrcorner v=(\omega\lrcorner v)\wedge\eta+(-1)^k\omega\wedge(\eta\lrcorner v),\\
\omega\in\operatorname{Alt}^kV,\ \eta\in\operatorname{Alt}^lV.
\end{gathered}
\]

% block: alternating-inner-product page:75 type:ordinary-paragraph
\medskip\noindent\textbf{\BT{Inner product}{内积}}\par
\BT{If the vector space $V$ is endowed with an \emph{inner product} $\langle\cdot,\cdot\rangle$, then this induces the Riesz isomorphism $V$ to $V'$ and so an inner product on $V'$. An \emph{inner product} on $\operatorname{Alt}^kV$ is then defined by}{若向量空间 $V$ 赋有\emph{内积} $\langle\cdot,\cdot\rangle$, 就诱导了从 $V$ 到 $V'$ 的 Riesz 同构, 从而在 $V'$ 上诱导内积. 于是可在 $\operatorname{Alt}^kV$ 上定义\emph{内积}}
\[
\langle u^1\wedge\cdots\wedge u^k,v^1\wedge\cdots\wedge v^k\rangle=\det\big[(\langle u^i,v^j\rangle)_{i,j=1}^n\big],
\]
\BT{which extends to general elements of $\operatorname{Alt}^kV$ by bilinearity. If the $u_i$ form an \emph{orthonormal} basis for $V$, then the corresponding dual basis elements $u^i$ are orthonormal in $V'$, as are the resulting basis functions \eqref{eq:alternating-form-basis} of $\operatorname{Alt}^kV$.}{并通过双线性性延拓到 $\operatorname{Alt}^kV$ 的一般元素. 若 $u_i$ 构成 $V$ 的\emph{标准正交}基, 则相应对偶基元素 $u^i$ 在 $V'$ 中标准正交, 由此得到的 $\operatorname{Alt}^kV$ 的基函数 \eqref{eq:alternating-form-basis} 也标准正交.}

% block: volume-form page:76 type:ordinary-paragraph
\medskip\noindent\textbf{\BT{Volume form}{体积形式}}\par
\BT{When $k$ is equal to the dimension $n$ of $V$, the space $\operatorname{Alt}^nV$ is one-dimensional, and so a basis consists of a single element. If the space $V$ is endowed with both an inner product and an \emph{orientation}, then there is a canonical choice of basis element. The ordered bases of a finite dimensional vector space $V$ divide into two classes according to whether the change-of-basis matrix connecting two of them has a positive or negative determinant. We orient the vector space by designating the ordered bases of one class as the positively oriented ones. (In $\mathbb R^3$, it is conventional to take as the positively oriented bases those which satisfy the "right-hand rule.") If $\omega\in\operatorname{Alt}^nV$, then the value of $\omega(u_1,\ldots,u_n)$ is the same for any orthonormal positively oriented basis $u_1,\ldots,u_n$, and we determine a unique element of $\operatorname{Alt}^nV$ by requiring this value to be $1$. The resulting alternating $n$-form is called the \emph{volume form} on $V$, which we typically denote by $\operatorname{vol}$. It may be expressed in terms of the dual basis to any positively oriented orthonormal basis as $\operatorname{vol}=u^1\wedge\cdots\wedge u^n$. On $\mathbb R^n$, the canonical basis $e_1,\ldots,e_n$ is orthonormal and positively oriented, so}{当 $k$ 等于 $V$ 的维数 $n$ 时, 空间 $\operatorname{Alt}^nV$ 是一维的, 因而一组基只包含一个元素. 若空间 $V$ 既赋有内积又赋有\emph{定向}, 就有典范的基元素选取. 有限维向量空间 $V$ 的有序基分为两类, 取决于连接其中两组基的换基矩阵的行列式为正还是为负. 指定其中一类的有序基为正向基, 就给向量空间定了向. (在 $\mathbb R^3$ 中, 习惯上将满足"右手法则"的基取为正向基.) 若 $\omega\in\operatorname{Alt}^nV$, 则对任意标准正交正向基 $u_1,\ldots,u_n$, $\omega(u_1,\ldots,u_n)$ 的值相同; 要求此值为 $1$, 就唯一确定了 $\operatorname{Alt}^nV$ 的一个元素. 得到的交错 $n$-形式称为 $V$ 上的\emph{体积形式}, 通常记为 $\operatorname{vol}$. 它可用任意正向标准正交基的对偶基写成 $\operatorname{vol}=u^1\wedge\cdots\wedge u^n$. 在 $\mathbb R^n$ 上, 典范基 $e_1,\ldots,e_n$ 标准正交且正向, 所以}
\[
\operatorname{vol}=dx^1\wedge\cdots\wedge dx^n.
\]

% block: hodge-star-definition page:76 type:ordinary-paragraph
\medskip\noindent\textbf{Hodge \BT{star}{星算子}}\par
\BT{We continue with an oriented inner product space $V$. Let $\omega\in\operatorname{Alt}^kV$ for some $0\leq k\leq n$, and let $\mu\in\operatorname{Alt}^{n-k}V$. Then $\omega\wedge\mu=c\operatorname{vol}$ for some real number $c$, and the map $\mu\mapsto c$ is a linear functional on $\operatorname{Alt}^{n-k}$. By the Riesz representation theorem it is given by the inner product of $\mu$ with some other element of $\operatorname{Alt}^{n-k}V$. That element depends on $\omega$ and is called the \emph{Hodge star} or Hodge dual of $\omega$, denoted $\star\omega$. The Hodge star is thus characterized by the relation}{继续考虑有定向的内积空间 $V$. 设 $\omega\in\operatorname{Alt}^kV$, 其中 $0\leq k\leq n$, 并设 $\mu\in\operatorname{Alt}^{n-k}V$. 则对某个实数 $c$, 有 $\omega\wedge\mu=c\operatorname{vol}$, 且映射 $\mu\mapsto c$ 是 $\operatorname{Alt}^{n-k}$ 上的线性泛函. 由 Riesz 表示定理, 它可表示为 $\mu$ 与 $\operatorname{Alt}^{n-k}V$ 中某个其他元素的内积. 该元素依赖于 $\omega$, 称为 $\omega$ 的 \emph{Hodge 星算子}或 Hodge 对偶, 记为 $\star\omega$. 因此 Hodge 星算子由下述关系刻画:}
\begin{equation}\label{eq:hodge-star-characterization}
\omega\wedge\mu=\langle\star\omega,\mu\rangle\operatorname{vol},\qquad\omega\in\operatorname{Alt}^kV,\ \mu\in\operatorname{Alt}^{n-k}V.
\end{equation}

% block: hodge-star-basis-and-square page:76 type:ordinary-paragraph
\BT{For each $k$, we thus have a linear map $\star:\operatorname{Alt}^kV\to\operatorname{Alt}^{n-k}V$. If we choose a positively oriented orthonormal basis of $1$-forms, $u^1,\ldots,u^n$, so $\operatorname{vol}=u^1\wedge\cdots\wedge u^n$, we find that}{因此, 对每个 $k$, 都有线性映射 $\star:\operatorname{Alt}^kV\to\operatorname{Alt}^{n-k}V$. 若选取 $1$-形式的正向标准正交基 $u^1,\ldots,u^n$, 使 $\operatorname{vol}=u^1\wedge\cdots\wedge u^n$, 则有}
\[
\star u^{\sigma_1}\wedge\cdots\wedge u^{\sigma_k}=\pm u^{\tau_1}\wedge\cdots\wedge u^{\tau_{n-k}}
\]
\BT{where $\sigma$ and $\tau$ are increasing sequences which partition $\{1,\ldots,n\}$ and the sign is that of the permutation $(\sigma,\tau)$. For example, in $\mathbb R^3$,}{其中 $\sigma$ 和 $\tau$ 是将 $\{1,\ldots,n\}$ 分割的递增序列, 符号是排列 $(\sigma,\tau)$ 的符号. 例如, 在 $\mathbb R^3$ 中,}
\[
\star dx^1=dx^2\wedge dx^3,\qquad\star dx^2=-dx^1\wedge dx^3,\qquad\star dx^3=dx^1\wedge dx^2.
\]
\BT{Note that $\star$ is an \emph{isometry} and $\star\star\omega=(-)^{k(n-k)}\omega$, or, equivalently, if $\omega\in\operatorname{Alt}^kV$, then}{注意, $\star$ 是\emph{等距映射}, 且 $\star\star\omega=(-)^{k(n-k)}\omega$; 或等价地, 若 $\omega\in\operatorname{Alt}^kV$, 则}
\[
\star\star\omega=\omega\text{ in odd dimensions},\qquad\star\star\omega=(-)^k\omega\text{ in even dimensions}.
\]
\BT{Combining this with \eqref{eq:hodge-star-characterization} gives an alternate characterization of the Hodge star:}{结合 \eqref{eq:hodge-star-characterization}, 得到 Hodge 星算子的另一种刻画:}
\begin{equation}\label{eq:hodge-star-inner-product}
\omega\wedge\star\nu=\langle\omega,\nu\rangle\operatorname{vol},\qquad\omega,\nu\in\operatorname{Alt}^kV.
\end{equation}

% block: scalar-vector-proxies pages:76,77 type:ordinary-paragraph
\medskip\noindent\textbf{\BT{Scalar and vector proxies}{标量代理和向量代理}}\par
\BT{If $V$ is an oriented inner product space of dimension $n$, then there is a natural isomorphism of each of the spaces $\operatorname{Alt}^0V$ and $\operatorname{Alt}^nV$ onto the scalar space $\mathbb R$, and of each of the spaces $\operatorname{Alt}^1V$ and $\operatorname{Alt}^{n-1}V$ onto the vector space $V$. In this way we can associate either a \emph{scalar proxy} or a \emph{vector proxy} to each member of $\operatorname{Alt}^kV$ for $k=0$, $1$, $n-1$, or $n$ (and so for all cases if $\dim V\leq3$). We now describe these isomorphisms.}{若 $V$ 是维数为 $n$ 的有定向内积空间, 则空间 $\operatorname{Alt}^0V$ 和 $\operatorname{Alt}^nV$ 各自与标量空间 $\mathbb R$ 之间存在自然同构, 空间 $\operatorname{Alt}^1V$ 和 $\operatorname{Alt}^{n-1}V$ 各自与向量空间 $V$ 之间也存在自然同构. 由此, 当 $k=0$, $1$, $n-1$ 或 $n$ 时 (因而在 $\dim V\leq3$ 时涵盖所有情形), 可为 $\operatorname{Alt}^kV$ 的每个元素关联一个\emph{标量代理}或\emph{向量代理}. 现在描述这些同构.}

% block: proxy-isomorphisms page:77 type:ordinary-paragraph
\BT{For any vector space $V$, $\operatorname{Alt}^0V$ is just the space $\mathbb R$ of scalars and the isomorphism is the identity. The inner product on $V$ defines the isomorphism of $\operatorname{Alt}^1V=V'$ onto $V$ via the Riesz representation. The Hodge star then maps $\operatorname{Alt}^{n-1}V$ isomorphically to $\operatorname{Alt}^1V$ (which is isomorphic to $V$) and similarly maps $\operatorname{Alt}^nV$ to $\operatorname{Alt}^0V=\mathbb R$.}{对任意向量空间 $V$, $\operatorname{Alt}^0V$ 就是标量空间 $\mathbb R$, 相应同构为恒等映射. $V$ 上的内积通过 Riesz 表示定义了从 $\operatorname{Alt}^1V=V'$ 到 $V$ 上的同构. 随后, Hodge 星算子将 $\operatorname{Alt}^{n-1}V$ 同构地映到 $\operatorname{Alt}^1V$ (后者与 $V$ 同构), 并类似地将 $\operatorname{Alt}^nV$ 映到 $\operatorname{Alt}^0V=\mathbb R$.}

% block: two-dimensional-proxies page:77 type:ordinary-paragraph
\BT{The case $n=2$ is particular, since in that case we have two ways of associating a vector proxy with an alternating $1$-form: either as a $1$-form by applying the Riesz map or as an $(n-1)$-form, that is, by applying first the Hodge star and then the Riesz map. These do not coincide, since the Hodge star $\operatorname{Alt}^1V\to\operatorname{Alt}^1V$ is not the identity map: it takes $u^1$ to $u^2$ and $u^2$ to $-u^1$. Therefore the two vector proxies differ by a rotation through a right angle.}{情形 $n=2$ 比较特殊, 因为此时有两种方式将向量代理与交错 $1$-形式关联: 或者将它视为 $1$-形式并应用 Riesz 映射, 或者视为 $(n-1)$-形式, 即先应用 Hodge 星算子再应用 Riesz 映射. 这两种方式并不一致, 因为 Hodge 星算子 $\operatorname{Alt}^1V\to\operatorname{Alt}^1V$ 不是恒等映射: 它将 $u^1$ 映为 $u^2$, 将 $u^2$ 映为 $-u^1$. 因此, 两个向量代理相差一个直角旋转.}

% block: linear-pullback page:77 type:ordinary-paragraph
\medskip\noindent\textbf{\BT{Pullback under linear maps}{线性映射下的拉回}}\par
\BT{If $L:V\to W$ is a linear map between finite dimensional vector spaces, then there is an obvious \emph{pullback} operation $L^*:\operatorname{Alt}^kW\to\operatorname{Alt}^kV$ given by}{若 $L:V\to W$ 是有限维向量空间之间的线性映射, 则有一个显然的\emph{拉回}运算 $L^*:\operatorname{Alt}^kW\to\operatorname{Alt}^kV$, 定义为}
\begin{equation}\label{eq:linear-pullback-definition}
L^*\omega(v_1,\ldots,v_k)=\omega(Lv_1,\ldots,Lv_k).
\end{equation}
\BT{The pullback acts contravariantly: if $K:W\to X$ is another map, then $(K\circ L)^*=L^*\circ K^*$. It also respects the algebraic structure, including the exterior product:}{拉回以反变方式作用: 若 $K:W\to X$ 是另一个映射, 则 $(K\circ L)^*=L^*\circ K^*$. 它也保持代数结构, 包括外积:}
\begin{equation}\label{eq:linear-pullback-wedge}
L^*(\omega\wedge\mu)=L^*\omega\wedge L^*\mu.
\end{equation}

% block: pullback-inclusion-and-normal page:77 type:ordinary-paragraph
\BT{In the case when $V$ is a subspace of $W$ and $L$ is the inclusion, the pullback of a $k$-form on $W$ is its restriction to $V\times\cdots\times V$. Assuming that $W$ has an inner product, we may pass to vector proxies to view the pullback $\operatorname{Alt}^1W\to\operatorname{Alt}^1V$ as a map from $W\to V$, which is nothing other than the orthogonal projection. Another example, is when $W$ is an oriented inner product space of dimension $n$ and $V$ an oriented subspace of dimension $n-1$. In that case, we may use proxies to view the pullback $\operatorname{Alt}^{n-1}W\to\operatorname{Alt}^{n-1}V$ as a map from $W\to\mathbb R$. Let us calculate this map. We start with a positive orthonormal basis $u_2,\ldots,u_n$ of $V$ and prepend to it the unit normal $u_1$ of $V$ in $W$ with the sign chosen so that $u_1,\ldots,u_n$ is positively oriented. Let $u^1,\ldots,u^n$ denote the dual basis of $W'$. The pullback of $u^1$ to $V'$ vanishes, while the pullbacks $v^2,\ldots,v^n$ of the remaining $u^i$ give the basis of $V'$ dual to $u_2,\ldots,u_n$. Since pullback respects the exterior product, it follows that}{当 $V$ 是 $W$ 的子空间且 $L$ 是包含映射时, $W$ 上一个 $k$-形式的拉回就是它在 $V\times\cdots\times V$ 上的限制. 假设 $W$ 具有内积, 则可通过向量代理将拉回 $\operatorname{Alt}^1W\to\operatorname{Alt}^1V$ 看作从 $W\to V$ 的映射, 它正是正交投影. 另一个例子是, $W$ 为维数 $n$ 的有定向内积空间, $V$ 为维数 $n-1$ 的有定向子空间. 此时可使用代理, 将拉回 $\operatorname{Alt}^{n-1}W\to\operatorname{Alt}^{n-1}V$ 看作从 $W\to\mathbb R$ 的映射. 来计算此映射. 先取 $V$ 的正向标准正交基 $u_2,\ldots,u_n$, 在前面添上 $V$ 在 $W$ 中的单位法向量 $u_1$, 并选取其符号使 $u_1,\ldots,u_n$ 为正向. 用 $u^1,\ldots,u^n$ 表示 $W'$ 的对偶基. $u^1$ 到 $V'$ 的拉回为零, 其余 $u^i$ 的拉回 $v^2,\ldots,v^n$ 则构成 $V'$ 中与 $u_2,\ldots,u_n$ 对偶的基. 由于拉回保持外积, 有}
\[
\begin{gathered}
L^*(u^2\wedge\cdots\wedge u^n)=v^2\wedge\cdots\wedge v^n=\operatorname{vol}_V,\\
L^*(u^1\wedge u^2\wedge\cdots\wedge\widehat{u^i}\wedge\cdots\wedge u^n)=0.
\end{gathered}
\]
\BT{Transferred to proxies, the pullback map $L^*$ thus becomes}{转用代理表示, 拉回映射 $L^*$ 因而成为}
\[
u_1\mapsto1,\qquad u_i\mapsto0,\qquad i=2,\ldots,n.
\]
\BT{This map is simply the inner product of a vector in $W$ with the oriented unit normal $u_1$. To summarize, in terms of proxies,}{此映射就是 $W$ 中向量与有定向单位法向量 $u_1$ 的内积. 总之, 用代理表示时,}
\Needspace{6\baselineskip}
\begin{itemize}
\item\BT{the pullback of a $1$-form to a subspace is the orthogonal projection, and}{一个 $1$-形式到子空间的拉回就是正交投影;}
\item\BT{the pullback of an $n-1$ form to a subspace of codimension $1$ is the inner product with the normal.}{一个 $n-1$ 形式到余维为 $1$ 的子空间的拉回, 就是与法向量的内积.}
\end{itemize}

% block: exterior-algebra-r3-table page:78 type:table
\begin{table}[htbp]\centering\footnotesize
\caption{\BT{Exterior algebra of $\mathbb R^3$ in terms of scalar and vector proxy fields.}{用标量和向量代理场表示的 $\mathbb R^3$ 外代数.}}\label{tab:exterior-algebra-r3}
\setlength{\tabcolsep}{4pt}
\begin{tabular}{l|l}\hline
\multicolumn{2}{c}{\BT{Proxies}{代理}}\\\hline
$\operatorname{Alt}^0\mathbb R^3=\mathbb R$&$c\leftrightarrow c$\\
$\operatorname{Alt}^1\mathbb R^3\xrightarrow{\cong}\mathbb R^3$&$u_1dx^1+u_2dx^2+u_3dx^3\leftrightarrow u$\\
$\operatorname{Alt}^2\mathbb R^3\xrightarrow{\cong}\mathbb R^3$&$u_1dx^2\wedge dx^3-u_2dx^1\wedge dx^3+u_3dx^1\wedge dx^2\leftrightarrow u$\\
$\operatorname{Alt}^3\mathbb R^3\xrightarrow{\cong}\mathbb R$&$c\leftrightarrow c\,dx^1\wedge dx^2\wedge dx^3$\\\hline
\multicolumn{2}{c}{\BT{Exterior product}{外积}}\\\hline
$\wedge:\operatorname{Alt}^1\mathbb R^3\times\operatorname{Alt}^1\mathbb R^3\to\operatorname{Alt}^2\mathbb R^3$&$\times:\mathbb R^3\times\mathbb R^3\to\mathbb R^3$\\
$\wedge:\operatorname{Alt}^1\mathbb R^3\times\operatorname{Alt}^2\mathbb R^3\to\operatorname{Alt}^3\mathbb R^3$&$\cdot:\mathbb R^3\times\mathbb R^3\to\mathbb R$\\\hline
\multicolumn{2}{c}{\BT{Contraction with a vector $v\in\mathbb R^3$}{与向量 $v\in\mathbb R^3$ 的缩并}}\\\hline
$\lrcorner v:\operatorname{Alt}^1\mathbb R^3\to\operatorname{Alt}^0\mathbb R^3$&$v\cdot:\mathbb R^3\to\mathbb R$\\
$\lrcorner v:\operatorname{Alt}^2\mathbb R^3\to\operatorname{Alt}^1\mathbb R^3$&$v\times:\mathbb R^3\to\mathbb R^3$\\
$\lrcorner v:\operatorname{Alt}^3\mathbb R^3\to\operatorname{Alt}^2\mathbb R^3$&$v:\mathbb R\to\mathbb R^3$\quad$(c\mapsto cv)$\\\hline
\multicolumn{2}{c}{\BT{Inner product}{内积}}\\\hline
\BT{Inner product on $\operatorname{Alt}^k\mathbb R^3$}{$\operatorname{Alt}^k\mathbb R^3$ 上的内积}&\BT{Euclidean inner product on $\mathbb R$ or $\mathbb R^3$}{$\mathbb R$ 或 $\mathbb R^3$ 上的 Euclidean 内积}\\\hline
\multicolumn{2}{c}{\BT{Volume form}{体积形式}}\\\hline
$\operatorname{vol}=dx^1\wedge dx^2\wedge dx^3$&$(v_1,v_2,v_3)\mapsto\det[v_1|v_2|v_3]$\\\hline
\multicolumn{2}{c}{Hodge \BT{star}{星算子}}\\\hline
$\star:\operatorname{Alt}^0\mathbb R^3\to\operatorname{Alt}^3\mathbb R^3$&$\operatorname{id}:\mathbb R\to\mathbb R$\\
$\star:\operatorname{Alt}^1\mathbb R^3\to\operatorname{Alt}^2\mathbb R^3$&$\operatorname{id}:\mathbb R^3\to\mathbb R^3$\\
$\star:\operatorname{Alt}^2\mathbb R^3\to\operatorname{Alt}^1\mathbb R^3$&$\operatorname{id}:\mathbb R^3\to\mathbb R^3$\\
$\star:\operatorname{Alt}^3\mathbb R^3\to\operatorname{Alt}^0\mathbb R^3$&$\operatorname{id}:\mathbb R\to\mathbb R$\\\hline
\multicolumn{2}{c}{\BT{Pullback by a linear map $L:\mathbb R^3\to\mathbb R^3$}{线性映射 $L:\mathbb R^3\to\mathbb R^3$ 的拉回}}\\\hline
$L^*:\operatorname{Alt}^0\mathbb R^3\to\operatorname{Alt}^0\mathbb R^3$&$\operatorname{id}:\mathbb R\to\mathbb R$\\
$L^*:\operatorname{Alt}^1\mathbb R^3\to\operatorname{Alt}^1\mathbb R^3$&$L^{\mathrm T}:\mathbb R^3\to\mathbb R^3$\\
$L^*:\operatorname{Alt}^2\mathbb R^3\to\operatorname{Alt}^2\mathbb R^3$&$(\det L)L^{-1}:\mathbb R^3\to\mathbb R^3$\\
$L^*:\operatorname{Alt}^3\mathbb R^3\to\operatorname{Alt}^3\mathbb R^3$&$(\det L):\mathbb R\to\mathbb R$\quad$(c\mapsto c\det L)$\\\hline
\end{tabular}
\end{table}

% block: exterior-algebra-r3-summary page:78 type:ordinary-paragraph
\medskip\noindent\textbf{\BT{Exterior algebra of $\mathbb R^3$}{$\mathbb R^3$ 的外代数}}\par
\BT{Since the spaces $\operatorname{Alt}^0\mathbb R^n$ and $\operatorname{Alt}^n\mathbb R^n$ may be identified with $\mathbb R$ and the spaces $\operatorname{Alt}^1\mathbb R^n$ and $\operatorname{Alt}^{n-1}\mathbb R^n$ with $\mathbb R^n$, in the case of $n=3$ (or fewer) dimensions, all the operations of exterior algebra may be viewed as operations on scalars or vectors. These operations are tabulated in Table~\ref{tab:exterior-algebra-r3}, which is adapted from \BookCite{11}. We remark that the matrix $(\det L)L^{-1}$ which shows up as the pullback of a $2$-form in the penultimate line of the table is the adjugate matrix of $L$, which is defined (as the transpose of the cofactor matrix) even when $L$ is singular.}{由于空间 $\operatorname{Alt}^0\mathbb R^n$ 和 $\operatorname{Alt}^n\mathbb R^n$ 可与 $\mathbb R$ 等同, 空间 $\operatorname{Alt}^1\mathbb R^n$ 和 $\operatorname{Alt}^{n-1}\mathbb R^n$ 可与 $\mathbb R^n$ 等同, 在 $n=3$ (或更低) 维时, 外代数的所有运算都可视为标量或向量上的运算. Table~\ref{tab:exterior-algebra-r3} 列出了这些运算, 该表改编自 \BookCite{11}. 注意, 表中倒数第二行作为 $2$-形式拉回出现的矩阵 $(\det L)L^{-1}$ 是 $L$ 的伴随矩阵, 即使 $L$ 奇异, 它也可定义为代数余子式矩阵的转置.}
